% ============================================================================
%  CHAPTER 4: CHALLENGES AND SOLUTIONS   (expanded, deep-dive version)
%  Replaces the entire Chapter 4 body in main.tex.
% ============================================================================

\chapter{Challenges and Solutions}
\label{chap:challenges}

Building FinITR-AI was, more than anything, an exercise in working out what the
language model should never be allowed to touch. The recurring discovery was
that the parts of the task a model finds easy --- producing fluent prose,
recognising the rough shape of a problem --- are not the parts that determine
whether the output is correct, and the parts that do determine correctness are
exactly where a model is least dependable. A second theme, present in almost
every design decision, was a deliberate asymmetry in how errors are weighed:
because a missed risk can cost a taxpayer real money in penalty and interest
while an over-cautious flag costs them only a moment's attention, the system is
tuned throughout to err toward false positives. This chapter sets out the major
obstacles we encountered and how each was resolved.
Table~\ref{tab:challenges} summarises them, and the discussion that follows
expands on the ones that shaped the architecture most.

\begin{table}[H]
\centering
\caption{Major Implementation Challenges and Their Solutions}
\label{tab:challenges}
\begin{tabularx}{\textwidth}{|X|X|}
\hline
\textbf{Challenge} & \textbf{Solution} \\
\hline
Language models compute Indian tax incorrectly --- wrong slab boundaries,
omitted cess, mis-applied 87A rebate and FY 2025-26 marginal relief --- and do
so with full confidence. & All arithmetic was removed from the model. A
deterministic CalculatorTool hard-codes the FY 2025-26 slabs, surcharge, cess,
rebate, and marginal relief. The ablation shows this lifts tax accuracy from
54.2\% to 94.0\%. \\
\hline
Models hallucinate disallowed claims --- most often an 80C or HRA deduction
under the New Regime --- which a non-expert cannot catch. & A Critic agent with
an NLI faithfulness verifier and a legal-corpus retriever blocks ungrounded
claims and forces the responsible agent to re-run. The hallucination rate fell
from 18.4\% (no critic) to 6.8\%. \\
\hline
Real documents carry noise: OCR rounding and small TDS differences make Form 16
disagree with the AIS by a few rupees. & The orchestrator treats a
$\leq$3\% Form-16-versus-AIS gap as noise and adopts the AIS salary as ground
truth, scaling basic salary proportionally, so the downstream computation stays
stable. \\
\hline
Aggressive schedule inference inflates recall but creates false positives;
conservative inference misses genuine cases. & Schedule inference was made
strictly signal-based --- a schedule is added only on concrete evidence (a bank
interest credit, an AIS SFT entry, a broker CSV). This holds schedule precision
at 96.2\% while accepting a known, documented recall gap. \\
\hline
Risk scoring on a single signal can under-rate a severe case --- undeclared
crypto seen only as a bank deposit was scored HIGH rather than CRITICAL. & An
additive, evidence-weighted risk model with documented escalation rules was
adopted, and the false-negative-averse banding ensures no genuine HIGH or
CRITICAL case is ever scored LOW, even where the exact band is arguable. \\
\hline
Bank-statement descriptions are terse, code-laden, and full of Hinglish vendor
names that defeat simple keyword rules. & A three-stage classifier ---
high-precision regex, then kNN over multilingual sentence embeddings, then an
LLM fallback --- normalises and labels each transaction, reaching 94.2\% macro
F1 with 100\% F1 on every income-critical class. \\
\hline
Running a 7B model locally is several seconds slower than calling a cloud API. &
The $\sim$6.1 s average latency was accepted as the price of privacy: the whole
pipeline runs offline, no document leaves the machine, and the model sits on the
explanation path rather than the arithmetic path. \\
\hline
\end{tabularx}
\end{table}

\section*{Discussion of the Decisive Challenges}
\addcontentsline{toc}{section}{Discussion of the Decisive Challenges}

The first two rows of the table carried the project. Forbidding the model from
doing arithmetic was the single most consequential architectural decision, and
it is the one that separates FinITR-AI's tax accuracy from the general models we
benchmarked against. The instinct, when a language model gets a sum wrong, is to
prompt it more carefully or to give it a better example. We found that instinct
to be a dead end for tax: the failures are not failures of instruction but of
the underlying mechanism, and no amount of prompting reliably teaches a network
to apply a marginal-relief taper or to remember that the cess compounds on tax
plus surcharge. Replacing the model's mental arithmetic with a few hundred lines
of deterministic code did what better prompting could not, and the ablation in
the next chapter quantifies the gap at nearly forty percentage points.

The Critic loop was the second pillar, and it solved a subtler problem than the
calculator. Grounding alone is not sufficient, because a model can retrieve the
correct rule and still state an incorrect conclusion in the same breath --- it
can quote the text of Section 80C accurately and then recommend claiming it
under a regime that disallows it. What was needed was not more evidence but the
authority to act on a verdict. Giving the Critic the power to veto a claim and
return the work to its author, rather than merely lowering a confidence score,
is what turned grounding from a passive feature into an active control. The
re-run is bounded at three iterations precisely because an unbounded
self-correction loop is its own hazard; in practice almost every case is clean
on the first review or resolved on the second.

The classifier challenge is worth dwelling on because it is where the messiness
of real data is most visible. An Indian bank statement does not say ``salary from
Infosys''; it says something like ``WDL TFR NEFT-SALARY-INFOSYS BPM
LTD-MAR25-UTR123456''. A purely rule-based approach drowns in these variants,
and a purely learned approach wastes the easy cases and is slow. The cascade was
the resolution: rules clear the unambiguous bulk instantly, sentence embeddings
absorb the noisy middle including Hinglish merchants that no rule anticipates,
and the model is consulted only for the residue. The result is both fast and
accurate where accuracy matters, since every label with a tax consequence is one
the cheap, reliable stages resolve.

The remaining rows reflect a quieter kind of engineering judgement: choosing,
consciously and with the trade-off in view, to give up a little recall for
precision, a little speed for privacy, and a little simplicity for robustness to
noise. None of these were forced by circumstance; each was a decision about what
kind of errors a tax tool should prefer to make. A tool that occasionally asks
the user to confirm an item it was unsure about is behaving correctly. A tool
that silently asserts a schedule the documents do not support, or that leaks a
salary figure to a remote server to save four seconds, is not. The error
philosophy described at the start of this chapter is, in the end, the thread
that ties these smaller decisions together.
